\section{Results}
\label{sec:results}

\subsection{Frozen reference-pocket redocking}

Table~\ref{tab:frozen_benchmark} summarizes the active frozen-manifest
evaluation. Pure confidence selects 488 poses below \(2\,\angstrom\) among
678 complexes, compared with 483 for Vina+DG and 486 for the historical
composite selector. The corresponding pure-confidence success rates are
76.47\% on Astex (65/85; 95\% Wilson interval 66.43--84.22\%), 73.05\% on
PoseBusters (225/308; 67.84--77.70\%), and 69.47\% on CASF (198/285;
63.90--74.53\%).

\begin{table}[H]
    \centering
    \caption{Top-1 symmetry-aware heavy-atom RMSD \(<2\,\angstrom\) under
    the frozen reference-pocket N80/S25 protocol. Oracle-80 is diagnostic
    only. The composite is retained for historical compatibility, whereas
    pure confidence is the deployable learned selector.}
    \label{tab:frozen_benchmark}
    \begingroup
    \small
    \setlength{\tabcolsep}{3.5pt}
    \begin{tabular}{lrrrrrr}
        \toprule
        Dataset & \(N\) & First (\%) & Vina+DG (\%) & Pure conf. (\%) &
        Frozen comp. (\%) & Oracle-80 (\%) \\
        \midrule
        Astex Diverse & 85  & 32.94 & 77.65 & 76.47 & 78.82 & 95.29 \\
        PoseBusters v2 & 308 & 35.71 & 71.10 & 73.05 & 72.73 & 94.81 \\
        CASF-2016 & 285 & 32.98 & 69.47 & 69.47 & 68.42 & 91.93 \\
        \bottomrule
    \end{tabular}
    \endgroup
\end{table}

The median RMSD of pure-confidence selections is \(1.156\,\angstrom\),
\(1.211\,\angstrom\), and \(1.274\,\angstrom\) on Astex, PoseBusters, and
CASF, respectively. Learned ranking improves over the same-candidate Vina+DG
baseline on PoseBusters but not uniformly across all three datasets. Likewise,
the frozen composite does not dominate pure confidence. These results support
the simpler minimum-predicted-RMSD rule as the deployment default but do not
establish superiority to external docking methods, whose pocket information,
candidate budgets, receptor preparation, and post-processing may differ.

\subsection{Candidate generation and pose selection are distinct bottlenecks}

Oracle-80 success exceeds 91\% on every dataset, leaving gaps of 18.82,
21.76, and 22.46 percentage points between pure-confidence top-1 and oracle
coverage on Astex, PoseBusters, and CASF. Thus most benchmark complexes have a
near-native candidate, but the ranker does not consistently identify it.
The gap is not attributed solely to confidence: the first-pose rates near
33--36\% show that candidate ordering from the stochastic generator is not
itself a useful selection rule.

The retained confidence checkpoint was selected on PLINDER validation rather
than on an external benchmark. On the frozen first 512 validation complexes,
minimum predicted RMSD gives 58.40\% top-1 success and the candidate oracle
gives 77.93\%. Four registered 1,500-step continuation studies---an unchanged
control, setwise, pairwise, and RMSD-listwise training---finish between 57.62
and 58.40\% and none passes the pre-registered \(+1.5\)-point gate. A separate
cluster-free filter changes only 3.72\% of selections on the full 1,076-complex
validation split and improves success from 53.90\% to 54.18\%, below its
\(1.0\)-point deployment gate. We therefore retain the unmodified step-42,500
confidence checkpoint and report the unsuccessful variants rather than
selecting post hoc on the external sets.

\subsection{Pocket context and center error}

Table~\ref{tab:pocket_sensitivity} reports pure-confidence top-1 success in
the separately frozen V2 sensitivity run. With an accurate center, a
\(6\,\angstrom\) crop is 25.88 points below \(10\,\angstrom\) on Astex and
28.57 points below it on PoseBusters. The best accurate-center crop is
\(8\,\angstrom\) for Astex and \(12\,\angstrom\) for PoseBusters, so no
single alternative consistently improves over the \(10\,\angstrom\) matched
default.

\begin{table}[H]
    \centering
    \caption{Pure-confidence top-1 RMSD \(<2\,\angstrom\) (\%) as a function
    of receptor crop radius and Gaussian center perturbation. Each row is one
    full-dataset condition of the promoted N80/S25 stack.}
    \label{tab:pocket_sensitivity}
    \begingroup
    \small
    \setlength{\tabcolsep}{8pt}
    \begin{tabular}{llrrrr}
        \toprule
        Dataset & Center jitter & \(6\,\angstrom\) & \(8\,\angstrom\) &
        \(10\,\angstrom\) & \(12\,\angstrom\) \\
        \midrule
        Astex (\(N=85\)) & \(0\,\angstrom\) & 50.59 & \textbf{77.65} & 76.47 & 76.47 \\
        & \(1\,\angstrom\) & 35.29 & \textbf{67.06} & \textbf{67.06} & \textbf{67.06} \\
        & \(2\,\angstrom\) & 29.41 & \textbf{45.88} & 42.35 & 40.00 \\
        \addlinespace
        PoseBusters (\(N=308\)) & \(0\,\angstrom\) & 44.16 & 66.56 & 72.73 & \textbf{74.03} \\
        & \(1\,\angstrom\) & 39.61 & 53.90 & \textbf{61.04} & 60.71 \\
        & \(2\,\angstrom\) & 20.78 & 31.17 & \textbf{34.42} & 31.17 \\
        \bottomrule
    \end{tabular}
    \endgroup
\end{table}

At the matched \(10\,\angstrom\) crop, \(1\,\angstrom\) center noise reduces
top-1 success from 76.47\% to 67.06\% on Astex and from 72.73\% to 61.04\%
on PoseBusters. At \(2\,\angstrom\), the rates fall to 42.35\% and 34.42\%.
Oracle-80 success at the same crop decreases from 94.12\% to 64.71\% on
Astex and from 93.18\% to 55.52\% on PoseBusters between zero and
\(2\,\angstrom\) perturbation. Center error therefore reduces generator
coverage, not only confidence ranking. Because this is an external-set
characterization, it does not by itself authorize changing the
\(10\,\angstrom\) inference default.

\subsection{Performance degrades for large and flexible ligands}

The same-candidate slices identify ligand size and flexibility as major failure
axes. On PoseBusters, pure-confidence success is 78.33\% for ligands with at
most 20 heavy atoms (\(N=120\)) but 44.00\% above 40 heavy atoms (\(N=25\));
oracle coverage decreases from 99.17\% to 60.00\%. For ligands with at least
eight rotatable bonds (\(N=66\)), pure-confidence and oracle success are
62.12\% and 86.36\%, compared with 77.88\% and 98.23\% for at most three
rotatable bonds (\(N=113\)). CASF shows the same size trend: pure-confidence
success drops from 82.57\% at at most 20 heavy atoms (\(N=109\)) to 38.89\%
above 40 atoms (\(N=18\)), with oracle coverage dropping from 97.25\% to
55.56\%. These paired declines indicate that larger ligands stress both
fragment-pose generation and candidate selection.

\subsection{Pose accuracy does not guarantee physical validity}

Official PoseBusters evaluation is currently available for the frozen
composite-selected poses, not for the pure-confidence selections. It gives a
pass-all rate of 54.87\% (169/308) when the RMSD criterion is excluded,
substantially below the corresponding composite RMSD success of 72.73\%.
Protein minimum-distance and tetrahedral-chirality checks pass for 69.16\% and
78.25\% of complexes, respectively, whereas internal clash, energy,
bond-length, and bond-angle checks each exceed 95\%. Near-native placement and
physical/chemical validity are therefore distinct failure axes.

A registered inference-only Vina-guidance study provides a negative result.
Guidance raises composite-selected official validity from 54.87\% to 56.17\%
without changing top-1 RMSD success, but the paired gain is only 1.30
percentage points (95\% bootstrap interval \(-0.32\) to \(3.25\) points;
exact McNemar \(p=0.289\)). It does not meet the pre-registered \(+3\)-point
validity target and is not included in the promoted default.
